\chapter{État de l'art et choix technologiques}
\markboth{Chapitre 2}{}

\section*{Introduction}

Le chapitre précédent a présenté le cadre du projet, analysé les solutions existantes
et défini les besoins fonctionnels et non fonctionnels des quatre modules.
Ce chapitre présente une étude des approches existantes dans les domaines couverts par le projet,
et justifie les choix technologiques retenus pour chacun des quatre modules.
Il est organisé en six axes : les assistants conversationnels,
la recherche documentaire intelligente, les modèles de génération de langage,
l'assistant émotionnel par intelligence artificielle,
la détection automatique du harcèlement,
et l'orchestration des pipelines conversationnels.
Chaque axe se conclut par une justification du choix retenu.

% ─────────────────────────────────────────────────────────────────────────────
\section{Assistants conversationnels}

\subsection{Évolution des approches}

Un assistant conversationnel est un programme capable de tenir une conversation
avec un utilisateur humain en langage naturel.
Ces systèmes ont évolué en plusieurs générations successives.

\textbf{Chatbots à règles.}
Les premiers systèmes reposaient sur des règles explicites du type
« si l'utilisateur dit X, répondre Y ».
Le programme fondateur est \textit{ELIZA}, créé en 1966 \cite{eliza1966}.
Ces systèmes sont rapides et prévisibles, mais incapables de traiter
des formulations non anticipées.

\textbf{Chatbots à intentions.}
À partir des années 2010, des plateformes comme Dialogflow \cite{dialogflow} ont introduit
la classification d'intention à partir d'exemples d'entraînement.
Ces systèmes sont plus souples, mais restent limités aux catégories définies
à l'avance et nécessitent un travail d'annotation important.

\textbf{Chatbots à grand modèle de langage.}
Les LLMs génèrent des réponses en langage naturel sans entraînement spécifique, mais peuvent produire des informations incorrectes — phénomène appelé \textit{hallucination}.

\textbf{Génération ancrée dans la documentation (RAG).}
L'architecture RAG (Retrieval-Augmented Generation — génération augmentée
par récupération) combine un LLM et une base de connaissances \cite{lewis2020rag}.
Le modèle répond uniquement à partir des documents fournis,
ce qui réduit significativement les hallucinations
et permet de contrôler le périmètre des réponses.

\subsection{Tableau comparatif}

Le Tableau~\ref{tab:comparatif_chatbots} compare les quatre approches
selon les critères pertinents pour le projet.

\vspace{6pt}
\begin{table}[H]
\centering
\caption{Comparaison des approches pour les assistants conversationnels}
\label{tab:comparatif_chatbots}
\begin{tabular}{|l|c|c|c|c|}
\hline
\textbf{Critère} & \textbf{Règles} & \textbf{Intentions} & \textbf{LLM seul} & \textbf{RAG} \\
\hline
Flexibilité & Faible & Moyenne & Élevée & Élevée \\
Fiabilité des réponses & Élevée & Moyenne & Faible & Élevée \\
Contrôle du périmètre & Élevé & Élevé & Faible & Élevé \\
Besoin de données & Aucun & Moyen & Aucun & Documents \\
Mise à jour facile & Non & Non & Non & Oui \\
\hline
\end{tabular}
\end{table}

\subsection{Choix retenu}

L'architecture RAG a été retenue pour les modules Noa assistant et assistant technique.
Elle combine la flexibilité d'un LLM avec la fiabilité d'une base documentaire maîtrisée,
ce qui est indispensable dans un contexte où les réponses incorrectes
peuvent avoir des conséquences sur des situations sensibles.
L'assistant émotionnel, en revanche, ne dispose pas de base documentaire fixe :
il repose sur un LLM en mode conversationnel pur, dont le comportement
est adapté à chaque échange selon le profil et l'état émotionnel de l'utilisateur.

% ─────────────────────────────────────────────────────────────────────────────
\section{Recherche documentaire intelligente}

\subsection{Approches existantes}

La recherche documentaire consiste à retrouver, parmi un ensemble de documents,
les passages les plus pertinents pour répondre à une question donnée.

\textbf{Recherche par mots-clés (TF-IDF).}
La méthode classique repose sur la fréquence des mots dans les documents.
Elle est rapide mais insensible au sens : une question formulée différemment
ne retrouvera pas les mêmes documents, même s'ils traitent du même sujet.

\textbf{Recherche par similarité sémantique.}
Les modèles d'encodage modernes transforment chaque texte
en un vecteur numérique \cite{sbert}.
Deux textes au sens similaire produisent des vecteurs proches,
mesurés par la similarité cosinus.
Cette approche est robuste aux reformulations et aux synonymes.

\textbf{Bases de données vectorielles.}
Les vecteurs générés sont indexés dans des bases spécialisées
qui permettent une recherche efficace parmi des milliers de documents.
Plusieurs solutions existent : Pinecone (hébergée dans le cloud),
FAISS (bibliothèque locale non persistante),
et ChromaDB (base locale persistante et légère).

\subsection{Choix retenus}

Le modèle d'encodage retenu est \texttt{paraphrase-multilingual-mpnet-base-v2} \cite{sbert},
pour sa compatibilité française et anglaise, sa robustesse aux reformulations
et son fonctionnement sans processeur graphique dédié.
La base vectorielle retenue est \textbf{ChromaDB} \cite{chromadb},
pour sa persistance locale, sa légèreté et son intégration aisée dans les pipelines Python.

% ─────────────────────────────────────────────────────────────────────────────
\section{Modèles de génération de langage}

\subsection{Panorama des modèles disponibles}

Plusieurs modèles de génération de langage sont disponibles pour des projets
nécessitant une inférence locale ou une maîtrise des données.

\textbf{GPT-4 (OpenAI).}
Modèle très performant, accessible via API.
Son utilisation implique que les données sont envoyées à des serveurs externes,
ce qui pose des problèmes de confidentialité dans le contexte scolaire.

\textbf{LLaMA (Meta).}
Famille de modèles ouverts utilisables localement.
Performants mais moins bien adaptés au français.

\textbf{Mistral (Mistral AI).}
Modèles développés par une entreprise française, disponibles en open source.
Le modèle \textbf{Mistral-7B-Instruct} \cite{mistral2023} est optimisé
pour suivre des instructions et offre un excellent rapport performance/ressources
pour la langue française.

\subsection{Tableau comparatif}

Le Tableau~\ref{tab:comparatif_llm} compare les principaux modèles
selon les critères du projet.

\vspace{6pt}
\begin{table}[H]
\centering
\caption{Comparaison des modèles de génération de langage}
\label{tab:comparatif_llm}
\begin{tabular}{|l|c|c|c|}
\hline
\textbf{Critère} & \textbf{GPT-4} & \textbf{LLaMA} & \textbf{Mistral-7B} \\
\hline
Hébergement local & Non & Oui & Oui \\
Qualité en français & Élevée & Moyenne & Élevée \\
Ressources requises & Cloud & Élevées & Modérées \\
Licence ouverte & Non & Oui & Oui \\
Confidentialité des données & Non & Oui & Oui \\
\hline
\end{tabular}
\end{table}

\subsection{Choix retenu}

Le modèle \textbf{Mistral-7B-Instruct} a été retenu pour sa qualité en français,
son hébergement local (garantissant la confidentialité des données des élèves),
et son accessibilité via le serveur LLM kaisens.fr.
La version AWQ, une version compressée du modèle, est utilisée
pour réduire les ressources mémoire sans perte significative de qualité.
Le serveur d'inférence \textbf{vLLM} \cite{vllm2023} permet de traiter
plusieurs requêtes simultanément.
Ce serveur est déjà déployé et opérationnel sur l'infrastructure de kaisens.fr,
ce qui élimine tout coût d'infrastructure supplémentaire et permet une intégration
immédiate sans mise en place d'un environnement GPU dédié.

% ─────────────────────────────────────────────────────────────────────────────
\section{assistant émotionnel par intelligence artificielle}

\subsection{Approches existantes}

L'assistant émotionnel automatisé vise à détecter les signaux de détresse
dans les messages d'un utilisateur et à générer des réponses empathiques adaptées.
Plusieurs approches ont été explorées dans la littérature et les applications commerciales.

\textbf{Scripts prédéfinis (chatbots guidés).}
Des applications comme Woebot ou Wysa reposent sur des arbres de décision
et des scripts thérapeutiques prédéfinis inspirés des techniques cognitivo-comportementales.
Ces systèmes sont prévisibles et sûrs, mais rigides :
ils ne peuvent pas s'adapter à une conversation libre
et offrent une expérience répétitive à long terme.

\textbf{Détection par mots-clés.}
Une approche simple consiste à rechercher des mots ou expressions
associés à une situation de crise (mal-être, idées suicidaires, violence).
Elle est rapide mais produit de nombreuses fausses alertes
et ne capture pas le contexte ni la nuance d'un message.

\textbf{Classifieurs d'apprentissage automatique.}
Des modèles supervisés entraînés sur des corpus annotés
permettent une détection plus fine de la détresse émotionnelle.
Ces modèles se concentrent sur la classification et ne génèrent pas de réponse empathique ;
ils nécessitent par ailleurs un corpus étiqueté conséquent, difficile à constituer.

\textbf{LLM conversationnel avec profil utilisateur.}
Un grand modèle de langage utilisé en mode conversationnel peut à la fois
détecter la détresse et générer une réponse empathique adaptée au contexte.
L'ajout d'un profil utilisateur persistant, mis à jour au fil des sessions,
permet au système de personnaliser les réponses en tenant compte
de la situation propre à chaque élève.

\subsection{Tableau comparatif}

Le Tableau~\ref{tab:comparatif_emotionnel} compare ces approches
selon les critères pertinents pour l'assistant émotionnel.

\vspace{6pt}
\begin{table}[H]
\centering
\caption{Comparaison des approches pour l'assistant émotionnel}
\label{tab:comparatif_emotionnel}
\begin{tabular}{|l|c|c|c|c|}
\hline
\textbf{Critère} & \textbf{Scripts} & \textbf{Mots-clés} & \textbf{ML supervisé} & \textbf{LLM conversationnel} \\
\hline
Flexibilité & Faible & Faible & Moyenne & Élevée \\
Qualité empathique & Moyenne & Nulle & Nulle & Élevée \\
Données d'entraînement & Aucune & Aucune & Élevées & Aucune \\
Personnalisation & Non & Non & Non & Oui \\
Détection de détresse & Manuelle & Partielle & Oui & Oui \\
\hline
\end{tabular}
\end{table}

\subsection{Choix retenu}

L'approche retenue est le \textbf{LLM conversationnel avec profil utilisateur persistant}.
Contrairement aux modules Noa assistant et assistant technique,
ce module n'a pas de base documentaire à consulter :
les réponses doivent s'adapter à l'état émotionnel de l'utilisateur,
pas à un corpus de documents.
Un profil est extrait automatiquement de la conversation et enrichi à chaque session,
permettant des réponses de plus en plus personnalisées.
Des garde-fous basés sur la classification du contenu détectent
les situations de crise et déclenchent une réponse adaptée
avec les numéros d'urgence appropriés.

% ─────────────────────────────────────────────────────────────────────────────
\section{Détection automatique du harcèlement}

\subsection{Approches existantes}

La détection automatique du harcèlement dans les textes est un problème
activement étudié en traitement automatique du langage naturel.

\textbf{Listes de mots offensants.}
L'approche la plus simple consiste à rechercher des mots ou expressions
préalablement identifiés comme offensants.
Elle est rapide mais incapable de saisir le contexte, l'ironie ou l'argot.

\textbf{Modèles d'apprentissage automatique classiques.}
Des modèles comme la régression logistique ou les SVM permettent
une classification binaire à partir de caractéristiques textuelles.
Ils nécessitent des données d'entraînement étiquetées et restent limités
pour les taxonomies complexes à plusieurs catégories.

\textbf{Modèles de type BERT.}
Les modèles basés sur l'architecture Transformer (comme CamemBERT pour le français)
offrent une meilleure compréhension du contexte.
Ils nécessitent un corpus d'entraînement étiqueté conséquent,
difficile à constituer dans le domaine du harcèlement scolaire en français.

\textbf{Génération ancrée dans des textes juridiques.}
Une approche originale consiste à guider le modèle de génération
par les textes de loi pertinents \cite{sparkloi}.
La décision est ainsi ancrée dans le droit français,
ce qui renforce la légitimité des résultats dans un contexte légal encadré.

\subsection{Tableau comparatif}

Le Tableau~\ref{tab:comparatif_classif} compare ces approches selon les critères du projet.

\vspace{6pt}
\begin{table}[H]
\centering
\caption{Comparaison des approches de détection du harcèlement}
\label{tab:comparatif_classif}
\begin{tabular}{|l|c|c|c|c|}
\hline
\textbf{Critère} & \textbf{Mots-clés} & \textbf{ML classique} & \textbf{BERT} & \textbf{RAG juridique} \\
\hline
Données d'entraînement & Aucune & Élevées & Élevées & Aucune \\
Compréhension du contexte & Faible & Moyenne & Élevée & Élevée \\
Taxonomie complexe & Non & Non & Partielle & Oui \\
Ancrage légal & Non & Non & Non & Oui \\
Mise à jour facile & Oui & Non & Non & Oui \\
\hline
\end{tabular}
\end{table}

\subsection{Choix retenu}

L'approche retenue est la \textbf{génération ancrée dans des textes juridiques},
combinée à un mécanisme de repli sur des règles heuristiques.
Cette combinaison offre à la fois la richesse sémantique d'un LLM
et la solidité d'un ancrage légal, sans nécessiter de corpus étiqueté.
Elle est implémentée en quatre stratégies comparées en parallèle,
détaillées dans le chapitre suivant.

% ─────────────────────────────────────────────────────────────────────────────
\section{Orchestration des pipelines conversationnels}

Les modules conversationnels du projet (Noa assistant, assistant technique, assistant émotionnel)
nécessitent une couche d'orchestration capable de chaîner plusieurs opérations :
classification d'intention, recherche documentaire, appel au modèle de langage
et gestion de la mémoire de session.
Plusieurs frameworks ont été étudiés pour remplir ce rôle.

\subsection{Frameworks disponibles}

\textbf{Implémentation personnalisée.}
Il est possible de gérer manuellement l'enchaînement des appels API
(ChromaDB, vLLM, Redis) sans framework dédié.
Cette approche offre un contrôle total mais impose un effort de développement
important pour chaque nouvelle fonctionnalité, sans bénéficier des abstractions
et optimisations déjà disponibles dans les outils spécialisés.

\textbf{LlamaIndex.}
LlamaIndex est un framework centré sur l'indexation et la recherche documentaire.
Il propose des connecteurs vers de nombreuses sources de données
et simplifie la construction de pipelines RAG.
Cependant, son périmètre est plus étroit que celui de LangChain :
il couvre moins bien la gestion de la mémoire conversationnelle,
les agents à boucle d'outils et l'orchestration de flux multi-étapes.

\textbf{LangChain.}
LangChain \cite{langchain} est un framework open source conçu spécifiquement
pour les applications fondées sur des modèles de langage.
Il fournit des abstractions pour les chaînes de traitement (\textit{chains}),
les agents (\textit{tool-use}), la mémoire conversationnelle, les connecteurs
vers les bases vectorielles et les modèles de langage.
Il supporte nativement ChromaDB, Redis et les API compatibles OpenAI,
ce qui correspond exactement à la pile technologique retenue pour ce projet.

\subsection{Tableau comparatif}

Le Tableau~\ref{tab:comparatif_orchestration} compare les trois approches
selon les besoins des modules conversationnels du projet.

\vspace{6pt}
\begin{table}[H]
\centering
\caption{Comparaison des frameworks d'orchestration}
\label{tab:comparatif_orchestration}
\begin{tabular}{|l|c|c|c|}
\hline
\textbf{Critère} & \textbf{Implémentation} & \textbf{LlamaIndex} & \textbf{LangChain} \\
                 & \textbf{personnalisée}  &                     &                    \\
\hline
Gestion mémoire conversationnelle & Non & Partielle & Oui \\
Support agents (tool-use)         & Non & Partiel   & Oui \\
Connecteur ChromaDB natif         & Non & Oui       & Oui \\
Connecteur Redis natif            & Non & Non       & Oui \\
Effort de développement           & Élevé & Moyen   & Faible \\
\hline
\end{tabular}
\end{table}

\subsection{Choix retenu}

LangChain a été retenu comme framework d'orchestration pour les trois modules conversationnels.
Il offre des connecteurs natifs vers ChromaDB et Redis, une gestion intégrée
de la mémoire de session et des abstractions pour les agents à boucle d'appels d'outils,
ce qui réduit significativement le volume de code à maintenir.
Sa compatibilité avec les API de style OpenAI permet de l'interfacer directement
avec le serveur vLLM de kaisens.fr sans adaptation supplémentaire.

% ─────────────────────────────────────────────────────────────────────────────
\section{Synthèse des choix technologiques}

Le Tableau~\ref{tab:synthese_choix} récapitule l'ensemble des choix technologiques retenus
et leur justification principale.

\vspace{6pt}
\begin{table}[H]
\centering
\caption{Synthèse des choix technologiques retenus}
\label{tab:synthese_choix}
\small
\setlength{\tabcolsep}{4pt}
\begin{tabular}{|p{4.5cm}|p{4cm}|p{6cm}|}
\hline
\textbf{Brique} & \textbf{Solution retenue} & \textbf{Justification principale} \\
\hline
Architecture chatbot (Noa assistant, assistant technique) & RAG & Fiabilité, contrôle du périmètre \\
Modèle d'encodage & MiniLM-L12-v2 & Bilingue, léger, sans GPU \\
Base documentaire & ChromaDB & Légère, locale, persistante \\
Modèle de génération & Mistral-7B-Instruct & Français, open source, local \\
Serveur d'inférence & vLLM & Requêtes simultanées \\
Orchestration & LangChain & Chaînes de traitement modulaires \\
Mémoire de session & Redis & Rapide, persistant \\
assistant émotionnel & LLM conversationnel & Flexibilité, empathie, pas de corpus \\
Profil utilisateur & Extraction LLM + Redis & Personnalisation multi-sessions \\
Déploiement & Docker & Portabilité, isolation \\
Classification texte & RAG juridique + repli & Ancrage légal, sans corpus étiqueté \\
Traitement en volume & Apache Spark & Traitement distribué et parallèle \\
\hline
\end{tabular}
\end{table}

% ─────────────────────────────────────────────────────────────────────────────
\section{Pile logicielle retenue}

Le Tableau~\ref{tab:env_logiciel} présente la pile logicielle complète du système.
Les versions et les rôles de chaque composant découlent directement des choix
justifiés dans ce chapitre.

\vspace{6pt}
\begin{table}[H]
\caption{Environnement logiciel global}
\label{tab:env_logiciel}
\centering
\small
\setlength{\tabcolsep}{3pt}
\begin{tabular}{|p{2.5cm}|p{4.5cm}|p{1.8cm}|p{5.5cm}|}
\hline
\textbf{Couche} & \textbf{Technologie} & \textbf{Version} & \textbf{Rôle} \\
\hline
Backend Python  & Python               & 3.11             & Langage principal — modules Noa assistant et assistant technique \\
\hline
Backend Python  & FastAPI \cite{fastapi} & $\geq$ 0.110   & Framework web — expose les interfaces de programmation REST \\
\hline
Backend Scala   & Scala                & 2.12.18          & Langage principal — module classificateur \\
\hline
Backend Scala   & Apache Spark         & 3.5.0            & Moteur de traitement de données à grande échelle \\
\hline
Orchestration   & LangChain \cite{langchain} & $\geq$ 0.2.0 & Orchestration des étapes de traitement conversationnel \\
\hline
Représentation sémantique & sentence-transformers \cite{sbert} & $\geq$ 2.7.0 & Calcul des vecteurs de représentation sémantique \\
\hline
Base documentaire & ChromaDB \cite{chromadb} & $\geq$ 0.5    & Stockage et recherche documentaire par similarité \\
\hline
Serveur LLM     & vLLM — kaisens.fr \cite{vllm2023} & —       & Serveur d'inférence externe, compatible API OpenAI \\
\hline
Modèle LLM      & Mistral-7B-Instruct \cite{mistral2023} & —   & Modèle de génération et de classification de texte \\
\hline
Sessions / Cache & Redis \cite{redis}  & $\geq$ 5.0       & Historique conversationnel et mémoire des réponses \\
\hline
Interface Noa assistant   & React + Vite         & 18 / 5           & Interface chatbot — site public \\
\hline
Interface assistant technique & Streamlit          & $\geq$ 1.32.0    & Interface assistant support interne \\
\hline
Déploiement     & Docker Compose \cite{docker} & —          & Conteneurisation et orchestration des services \\
\hline
\end{tabular}
\end{table}

% ─────────────────────────────────────────────────────────────────────────────
\section*{Conclusion}

Ce chapitre a présenté l'état de l'art et justifié les choix technologiques
retenus pour chaque module. L'ensemble des solutions sélectionnées forme une architecture
cohérente, adaptée aux contraintes de confidentialité et de déploiement du projet.
Le chapitre suivant présente la conception du système.
